\documentclass{article}
\usepackage[UTF8]{ctex}
\usepackage{amsmath,amssymb,amsfonts}
\usepackage{enumitem}
\usepackage{geometry}
\geometry{a4paper, margin=1in}
% === 错题本排版优化（不参与题目标识） ===
\setlength{\parindent}{0pt}
\setlength{\parskip}{0.35em plus 0.10em minus 0.08em}
\setlength{\emergencystretch}{3em}
\setlist{itemsep=0.25em, topsep=0.40em, parsep=0pt}
\allowdisplaybreaks[2]
% === 错题本排版优化结束 ===
\begin{document}
\section*{北京科技大学2024年高等代数真题}
\section*{9．北京科技大学2024年高等代数试卷}
一．（15 分）计算 \(n\) 阶行列式

\[
D_{n}=\left|\begin{array}{cccccc}
1 & 2 & 3 & \cdots & n-1 & n \\
2 & 3 & 4 & \cdots & n & 1 \\
3 & 4 & 5 & \cdots & 1 & 2 \\
\vdots & \vdots & \vdots & & \vdots & \vdots \\
n-1 & n & 1 & \cdots & n-3 & n-2 \\
n & 1 & 2 & \cdots & n-2 & n-1
\end{array}\right|
\]

二．（15 分）设 \(\alpha_{1}=(1,2,0)^{\mathrm{T}}, \alpha_{2}=(1, a+2,-3 a)^{\mathrm{T}}, \alpha_{3}=(-1,-b-2, \bar{a}+2 b)^{\mathrm{T}}, \beta=(1,3,-3)^{\mathrm{T}}\) ．试讨论当 \(a, b\) 为何值时。

\begin{itemize}

  \item[1．] \(\beta\) 不能由 \(\alpha_{1}, \alpha_{2}, \alpha_{3}\) 线性表出。

  \item[2．] \(\beta\) 可由 \(\alpha_{1}, \alpha_{2}, \alpha_{3}\) 唯一的线性表示，并求出表示式。

  \item[3．] \(\beta\) 可由 \(\alpha_{1}, \alpha_{2}, \alpha_{3}\) 线性表示，但表示式不唯一，并求出其一般表示式。
\end{itemize}

三．（15 分）设 \(\alpha\) 为3维单位列向量，\(E\) 为3阶单位矩阵，求矩阵 \(E-\alpha \alpha^{\mathrm{T}}\) 的秩．

四．（15 分）已知实对称矩阵

\[
A=\left(\begin{array}{ccc}
2 & 2 & -2 \\
2 & 5 & -4 \\
-2 & -4 & 5
\end{array}\right) .
\]

1．若 \(t E+A\) 为正定矩阵，求 \(t\) 的取值范围．\\

2．求正交矩阵 \(Q\) ，使得 \(Q^{\mathrm{T}} A Q\) 为对角矩阵。

\begin{itemize}
  \item[五．] （20分）设 \(A\) 为 \(m \times n\) 实矩阵，令
\[
N(A)=\left\{X \in \mathbb{R}^{n} \mid A X=0\right\},
\]
\(R(A)=\left\{Y \in \mathbb{R}^{m} \mid\right.\) 存在 \(X \in \mathbb{R}^{n}\) 使得 \(\left.A X=Y\right\}\).

  \item[1．] \(N(A)\) 为 \(\mathbb{R}^{n}\) 的子空间，\(R(A)\) 为 \(\mathbb{R}^{m}\) 的子空间。

  \item[2．] \(N(A)^{\perp}=R\left(A^{\mathrm{T}}\right)\) ，这里 \(N(A)^{\perp}\) 表示 \(N(A)\) 的正交补空间。
\end{itemize}

六。（20分）设 \(A, B, C, D\) 是数域 \(P\) 上的两两可交换的 \(n\) 阶方阵，且 \(A C+B D=E_{n}\) ，这里 \(E_{n}\) 为\(n\) 阶单位矩阵，记线性方程组 \(A B X=0, A X=0\) 和 \(B X=0\) 的解空间分别为 \(W, W_{1}, W_{2}\) 。证明：\(W=W_{1} \oplus W_{2}\) ．

七．（25 分）设 \(D=\left(\begin{array}{cc}A & C \\ C^{\mathrm{T}} & B\end{array}\right)\) 为正定矩阵，其中 \(A, B\) 分别为 \(m\) 阶和 \(n\) 阶对称矩阵，\(C\) 为 \(m \times n\) 矩阵，证明：\(B-C^{\mathrm{T}} A^{-1} C\) 为正定矩阵。

八．（10 分）设 \(p\) 为素数，\(a\) 为整数，且 \(p^{2} \mid(a+1)\) ，判断多项式 \(f(x)=a x^{p}+p x+1\) 在有理数域上是否可约，并给出理由。

九。（15 分）设 \(V\) 为复数域上的4维线性空间，\(e_{1}, e_{2}, e_{3}, e_{4}\) 为 \(V\) 的一组基，令 \(\varphi\) 为 \(V\) 上的线性变换，且 \(\varphi\) 在基 \(e_{1}, e_{2}, e_{3}, e_{4}\) 下的矩阵为

\[
A=\left(\begin{array}{cccc}
3 & -1 & 1 & -7 \\
-9 & -3 & -7 & -1 \\
0 & 0 & 4 & -8 \\
0 & 0 & 2 & -4
\end{array}\right)
\]

1．求 \(V\) 的一组基 \(f_{1}, f_{2}, f_{3}, f_{4}\) ，使得 \(\varphi\) 在该组基下的矩阵为若尔当标准形矩阵。\\

2 ．求从基 \(e_{1}, e_{2}, e_{3}, e_{4}\) 到基 \(f_{1}, f_{2}, f_{3}, f_{4}\) 的过渡矩阵．
\end{document}
